\documentclass[11pt,a4paper]{scrartcl}

%%%%%%%%%%%%%%%%%%%%%%%%%%%%%%%%%%%%%%%%%
%  Default packages
%%%%%%%%%%%%%%%%%%%%%%%%%%%%%%%%%%%%%%%%%
\usepackage{textcomp}%\textmu 
\usepackage{graphicx}
\usepackage{subfig}%allow subfigures
\usepackage{amsmath}%improve layout of formulas
\usepackage{color}
\usepackage{units}%\unitfrac (nice m/s in text and math mode)
\usepackage{url}
\usepackage{booktabs}%\toprule, \midrule, bottomrule in tables
\usepackage{xspace}%automatically add space at the end of macros
\usepackage{upgreek}%\upright greek letters in math environment
\usepackage{authblk}%several authors with different affiliations
\usepackage{threeparttable}% for footnotes within tables
\usepackage{multicol}%multimple columns

%% clickable links in pdf
\usepackage[pdftex,bookmarks=true,breaklinks=true,bookmarksnumbered=true,
            colorlinks=true,linkcolor=webgreen,citecolor=webred,urlcolor=webblue]{hyperref}
%% link colors
\definecolor{webred}{rgb}{0.5, 0, 0}
\definecolor{webgreen}{rgb}{0, 0.5, 0}
\definecolor{webblue}{rgb}{0, 0, 0.5}

%Command to allow footnotes within floats (tables, figures....)
%Use \floatingfootnotemark inside the float, which only places the mark and
%counts up the footnote counter.
%Immediately after the float you put the footnote text. We will see if it
%shows up on the right page....
\newcommand{\floatingfootnotemark}{\addtocounter{footnote}{1} \footnotemark[\value{footnote}]}
\newcommand{\floatingfootnotetext}[1]{\footnotetext[\value{footnote}]{#1}}

%A \fixme command with counter. At the end of the document you can display the counter to see how many FIXMEs are left
\newcounter{nFixmes}
\setcounter{nFixmes}{0}
\newcommand{\fixme}[1]{\addtocounter{nFixmes}{1}\textcolor{red}{(\includegraphics[height=2ex]{fixme} #1)}\xspace}

%file with the svn revision, generated using svnversion
\input{svnrevision.tex}

\title{Requirement for the MTCA4U Process Variable Adapter}
\date{\svnrevision\\ PDF build \today}
\author[1]{Martin Killenberg} \author[2]{Sebastian Marsching}
\author[3]{Adam Poitrokwsi}
\author[1]{Christian Schmidt}

\affil[1]{Deutsches Elektronen-Synchrotron DESY, Hamburg, Germany}
\affil[2]{aquenos GmbH, Baden-Baden, Germany}
\affil[3]{FastLogic Sp.\ z o.\ o., \L\'od\'z, Poland}

%%%%%%%%%%%%%%%%%%%%%%%%%%%%%%%%%%%%%%%%%
% This is where the document really begins
%%%%%%%%%%%%%%%%%%%%%%%%%%%%%%%%%%%%%%%%%
\begin{document}

\maketitle
\section{Objective and Scope}

MTCA4U's control system interface adapter is intended to decouple the business
logic of a server from the actual control system in use. For this it provides
two functionalities: 
\begin{enumerate}
  \item Control System Variable adapters
  \item Function callbacks to react on control system interrupts
\end{enumerate}

\subsection{Control System Variable adapters}

The control-system variable adapter is a thin layer which provides a
standardised interface for accessing the variable content, but does not provide
control-system functionality. In other words, a control system variable that
holds a signed 32 bit integer can be accessed as 32 bit integer from the
business logic, but a variable history, which might be provided by the control
system, is not accessible. The history is, however, accessible via the control
system. Like this, no control system functionality has to be abstracted or
duplicated/implemented if not available on certain control systems. 

The only exception is a list of callback functions which can be registered as
actions to be taken when the variable is changed. This is needed to allow the
synchronisation with the control system. The business logic can choose between
a set()-function which only informs the control system, or a set()-function
which also triggers the callbacks. When the variable is changed by the control
system, the callbacks are always triggered. A more detailed explanation of the
variable adapter can be found in
section~\ref{section_process_variable_adapter}. 

\subsection{Control System Callback Function}

Control systems can trigger actions on signals, user interrupts or might
implement events which represent an action which is not coupled to a process
variable. The control-system adapter implements an interface so that the
business logic can define functions to be executed, and the control-system
dependent code implements how the particular function is triggered by the
specific control system. 

\subsection{Python interface}

For easy interactive testing, it would be nice to have a python interface to
the control system stub. Like this one can run tests of the full business
logic without the necessity to install a control system. The interface might
also be used to implement tests. 

\section{Data Types}

\subsection{Simple Data Types}

The control system interface adapter supports ProcessVariable interface
adapters for the following simple data types\\[2ex] 
\begin{tabular}{rrrrr}
float & int8 & int16 & int32 & int64\\
double & uint8 & uint16 & uint32 & uint64 \\
string\\[2ex]
\end{tabular}

It is assumed that float, double, 16 bit integer and 32 bit integer are
supported by all control systems \fixme{Check at least with DOOCS and
  EPICS}. In case the short integers are not supported, the next larger
integer with range limitation is used. If no range limitations are available
in the control system, the maximum allowed integer for the shorter type is set
and reported back to the control system. In case 64 bit integers are not
supported two 32 bit integers are used. \fixme{How to encode signed int64?
  Alternative solution: MTCA4U does not support 64 bit integers} 
Control system specific data types are not supported.

\subsection{Strings}

Depending on the control system implementation the string might have limited
length. If an input string is too long it will be truncated. The string has a
“truncated” flag which can be queried. It intentionally does not throw an
exception in order not to avoid making servers crash which run fine on other
control systems. In most cases, there will be a truncated status message, which
can be tolerated. There might be, however, malfunctions if the complete string
is parsed automatically. In this case the parsing code will generate the error
rather than the server terminating with an exception. 

\subsection{Arrays}

MTCA4U supports arrays of all supported simple data types. \fixme{How to implement 64 bit integers if not supported by the CS?}

\section{Requirements for the ProcessVariable Adapter}\label{section_process_variable_adapter}

The ProcessVariable class holds an instance of the control system variable of
the corresponding type. The control system usually implements a thread and
locking scheme for the communication. This is the reason why ProcessVariables
cannot be considered real time capable. 

The business logic, however, will probably have a real time thread. For this
reason it cannot be avoided to do one copy: From the lock-free variable used
in the business logic into the control system variable. To avoid additional
copying the ProcessVariable adapter does not contain an instance of the basic
data type (or array), but accessors directly work with the control system
variable. \fixme{Should it be recommended that the business logic always use a
  copy of the variable and not directly use adaptor? It might be that the
  variable changes during a calculation if used multiple time, and it might be
  slow because of frequent locking/unlocking.}  

The business logic has to create instances of process variables without
knowing which particular implementation is needed. For this a factory pattern
is used. Each variable is identified by a unique name, which is used to
request creation of the variable from the factory. The actual address, under
which the variable is seen from the control system, is control system specific
and handled in the concrete factory. It is determined from the unique name. 

The ProcessVariable adapter is not thread safe. If the underlying control
system variables are thread safe, then the adapter is also thread safe as it
does no further caching but directly accesses the underlying variable. But this
is not guaranteed by the interface and the user has to implement a thread safe
mechanism himself. MTCA4U will provide thread support as a separate tool of
the Control System Tools package, probably implementing an actor pattern.
 \fixme{I want to prevent reinventing the
  wheel. Should we just use the actor pattern implemented in boost::asio? MK} 


\subsection{The ProcessVariable adapter interface features}
\begin{itemize}
\item a get() method
\item a set() method which triggers the callback functions
\item a list of callback functions which are executed if the variable contents
  changed
\item a set() method which does not trigger the callback functions. Like this
  it is possible the execute the callback functions only when the variable
  content is changed externally (from the control system), but update the
  variable from inside the business logic to inform the control system about
  the new value without triggering the callbacks. 
  \fixme{Do we want a second list of callbacks for this?}
\item an assignment operator (for convenience, like set operator with
  callback) \fixme{without callback?}
  \fixme{This will probably only possible with an additional copy} 
\item a const implicit conversion operator \fixme{Do we really need this?
  Business logic will use internal variables only. But how about simple,
  synchronous I/O servers?} 
\item \fixme{Do we want range limits and engineering units?}
\end{itemize}

\subsection{Additional feature of the array interfaces}
The array has an interface compatible with a boost::array/std::array.
\begin{itemize}
\item {[ ]} operator to access individual elements
\item iterators
\item size()
\item front()
\item back()
\item empty() \fixme{Do we need it ?}
\item swap() \fixme{Do we need it ?}
\item fill() \fixme{Do we need it ?}
\end{itemize}

\subsection{Implementation details}

\begin{itemize}
\item Process variables are not copy-constructible as they have to be
  registered with a unique name. They have a private copy constructor. 
\item The constructor is also private. Construction is done with a
  factory. The construction call takes a programme-wide unique name to
  identify the variable. With this the factory will be able to determine the
  required information to create the control system variable from a config
  file or data base.\\
N.B. Making the constructor private disables the possibility that the user
writes control-system dependent code. 
\fixme{This require the factory to be friend in order to access the
  constructor. Not so nice, but I don't see away around it which does not
  allow the user to write control system dependent code} 
\item \fixme{Do we allow arrays to change size?}
\item Call back functions can be unregistered.
\item In case signed variables have to be implemented as unsigned in the
  control system, or vice versa, a reinterpret-cast is done in order not to
  lose information. In this case the negative range/ upper half of the dynamic
  range has to be reinterpreted by the user. An 'out of range' flag is set. 
\item If only single precision floating point numbers are provided by the
  control system, the double is converted loosing precision. A 'limited
  precision' flag is set (statically in the concrete implementation.)
\end{itemize}

\subsection{Large data objects}

Some controls systems use synchronous operations like Remote Procedure Calls
for data transmission. This is not efficient for large data objects or objects
which have a very high update rate. In some cases (DOOCS for instance) the
control system has a special mechanism for these data, other control systems
like EPICS 3 usually require an 'out of band' implementation which adds a
proprietary protocol. Some control systems, however, use asynchronous data
protocols and can handle large data objects and high update rates efficiently
(EPICS 4, Tango, Karabo, Tine(?) ).

As the decision if a special channel is needed and from which size depends on
the control system, this has to be handled in the control system specific
ProcessVariable factory. For the business logic and in the interface the
notion of a special channel for high data rates does not exist.

\subsection{The ProcessVariable Factory}

The factory allows the abstraction of the variable creation. Each user class
which is programmed as a control system independent server is handed a
reference or shared pointer of a concrete factory implementation. The factory
is not implemented as a singleton pattern to allow the use of multiple
factories simultaneously, for instance one for low bandwidth channels and one
for high bandwidth channels (DAQ). This would allow to use the same 'out of
band' factory for multiple control systems with slow, synchronous calls. The
factory which is handed to the business logic would be a composite of the low
and high bandwidth factory in this case. 
For asynchronous control systems with sufficient bandwidth in the protocol the
factory would just have one component. All this is transparent to the business
logic.

The decision which channels is treated out of band is done inside the concrete
factory. It might be that one control system has to tread a certain amount of
data out of band while another control system can do this with its internal
protocol, but needs out of band for a higher data rate. Or that system A can
handle high update rates well, but has to treat large data blocks differently,
while system B can cope with the large data, but not with the high update
rate.

The factory is automatically instantiating and registering all variables
requested by the business logic. The server/application specific code should
not have to deal with the individual variables. It should only instantiate the
correct factory and pass it to the instance of the business logic class.

\subsection{Control System Stub}
For the unit testing a factory for a control system stub is implemented.
Stub variables can be set and modified via the ProcessVariable adapter and from
the C++ application which hosts the stub, incl. the callback functions on
change, but no control system functionality like network communication or
history is available. This allows to implement to run the business logic
stand-alone, enables full testing and even the possibility to write a GUI.

Nice to have: The control system stub should also have Python bindings so PyQT
or simple control scripts can be used, e.g. \fixme{incomplete sentence}

\section{Requirements for the Control System Function Callbacks}
\label{section_function_callbacks}
\fixme{to be written}

\section{Further, General Requirements}
As we have to support Ubuntu 12.4, C++11 is not used because the C++11 support
is not mature enough in this distribution. Boost is used for the required
functionality like shared pointers etc.

\appendix
\section{DOOCS Variable Types}

Directly supported by the control system adapter in bold
\begin{multicols}{3}
\noindent D\_name \\
\textbf{D\_float}\\
D\_float\_ptr \\
\textbf{D\_floatarray}\\
\textbf{D\_double}\\
D\_double\_ptr \\
\textbf{D\_doublearray}\\
D\_gen\_sts \\
D\_sts \\
D\_status \\
D\_polynom \\
D\_hist \\
D\_histStat \\
D\_comment \\
\textbf{D\_string} (80 characters version)\\
D\_zmqstring \\
D\_ustr \\
D\_text \\
D\_xml \\
D\_alarm \\
D\_devinfo \\
D\_plotinfo \\
D\_plotinfo\_p \\
D\_spectrum \\
D\_ifff \\
D\_filter \\
D\_dig\_filter \\
D\_bit \\
D\_xyzs \\
D\_iiii \\
\textbf{D\_int}\\
D\_int\_ptr \\
\textbf{D\_shortarray}\\
\textbf{D\_intarray}\\
\textbf{D\_longarray}\\
D\_call \\
D\_alarm \\
D\_exec \\
D\_error \\
D\_xy \\
D\_floatmean/D\_doublemean \\
D\_floatrms/D\_doublerms \\
D\_floatsigma/D\_doublesigma \\
D\_floatema/D\_doubleema \\
\textbf{D\_byte\_array}\\
D\_mdfloat\_array \\
D\_image \\
D\_imagec
\end{multicols}

Not supported by DOOCS:
\begin{itemize}
\item int8, int16 $\rightarrow$ use int32
\item int64 $\rightarrow$ \fixme{What to use?}
\item uint (all sizes) $\rightarrow$ use signed int with range check
\end{itemize}
The same for arrays.

\fixme{What is the difference between array and spectrum, except for some
  plotting meta information? How do we get a spectrum as currently used out of
  an array in the automated DOOCS dependent code (decide when to use array,
  when to use spectrum)?}

\section{EPICS Variable Types}
Directly supported by the control system adapter in bold
\begin{multicols}{3}
\noindent \textbf{DBF\_STRING}\\
\textbf{DBF\_INT}\\
\textbf{DBF\_SHORT}\\
\textbf{DBF\_FLOAT}\\
DBF\_ENUM\\
\textbf{DBF\_CHAR}\\
\textbf{DBF\_LONG}\\
\textbf{DBF\_DOUBLE}
\end{multicols}

Unsigned integer types are not fully supported, therefore the corresponding
signed type should be used and interpreted as unsigned. DBF\_STRING is
limited to 40 characters including the terminating null byte (so 39
characters effective size). Fixed-size arrays of all listed data-types
are also supported.

\section{FIXMEs}
Number of remaining FIXMEs: \arabic{nFixmes}

\end{document}
